\renewcommand{\Anno}{1996}

\renewcommand{\NumeroProblema}{2}

\renewcommand{\Testo}{
Si vuole circoscrivere ad un cerchio di raggio
$r$ un triangolo isoscele di perimetro $kr$.  Determinare per quali
valori del parametro $k$ il problema \`{e} risolubile.
}

\renewcommand{\Risposta}{$k\geq 6\sqrt{3}$}

\renewcommand{\Aiutino}{Occorre dimostrare due punti:
\begin{itemize}
	\item  che il perimetro minimo \`{e} quello del triangolo equilatero;

	\item che ogni perimetro superiore \`{e} raggiunto da qualche 
	triangolo isoscele.
\end{itemize}

La prima affermazione si riduce alla disuguaglianza $p\geq 
3\sqrt{3}\,r$ ($p$ indica il semiperimetro) e pu\`{o} essere dimostrata 
ricorrendo alla formula $r=S/p$, dove $S$ indica l'area, e applicando 
successivamente la formula di Erone per l'area e la disuguaglianza 
tra media geometrica e media aritmetica. Per la seconda affermazione 
basta mostrare che esistono triangoli isosceli con perimetro grande a 
piacere e concludere con un argomento di continuit\`{a}.}

\renewcommand{\SoluzioneUfficiale}{\TitoloSoluzione
Siano $a$, $b$, $c$ le lunghezze dei lati di un triangolo, sia 
$P=2p=a+b+c$ il suo perimetro, e sia $r$ il raggio della circonferenza 
inscritta. Dimostreremo che 
\begin{equation}
	P\geq 6\sqrt{3}\,r
	\label{eq:disug-Pr}
\end{equation}
con uguaglianza se e solo se il triangolo \`{e} equilatero, e che, per 
ogni coppia di numeri reali positivi $P$ ed $r$ soddisfacenti tale 
disuguaglianza, esiste un triangolo isoscele con perimetro $P$ e 
circonferenza inscritta di raggio $r$.

Questo baster\`{a} per concludere che il triangolo richiesto esiste se 
e solo se $k\geq 6\sqrt{3}$.

Per dimostrare la disuguaglianza (\ref{eq:disug-Pr}), ricordiamo che
\begin{equation}
	r=\frac{S}{p},
	\label{eq:r}
\end{equation}
dove $S$ indica l'area del triangolo, la quale per la formula di Erone 
\`{e} data da
\begin{equation}
	S=\sqrt{p(p-a)(p-b)(p-c)}.
	\label{eq:Erone}
\end{equation}
Dalla (\ref{eq:r}) e dalla (\ref{eq:Erone}) abbiamo pertanto che
\begin{equation}
	\frac{r}{2p}=\frac{S}{2p^{2}}=\frac{1}{2}
	\sqrt{\left(1-\frac{a}{p}\right)\left(1-\frac{b}{p}\right)
	\left(1-\frac{c}{p}\right)}.
	\label{eq:rsup}
\end{equation}

Ricordiamo ora la disuguaglianza tra media geometrica e media 
aritmetica
$$\sqrt[3]{xyz}\leq\frac{x+y+z}{3},$$
valida per ogni terna di numeri reali positivi $x$, $y$, $z$, con 
uguaglianza se e solo se $x=y=z$. Applicandola con
$$x=1-\frac{a}{p},
\hspace{3em}
y=1-\frac{b}{p},
\hspace{3em}
z=1-\frac{c}{p},$$
ed osservando che
$$x+y+z=3-\frac{a+b+c}{p}=3-\frac{2p}{p}=1,$$
otteniamo che 
$$\left(1-\frac{a}{p}\right)\left(1-\frac{b}{p}\right)
\left(1-\frac{c}{p}\right)\leq\frac{1}{27}.$$
	
Inserendo tale stima nella (\ref{eq:rsup}) troviamo infine che
$$\frac{r}{2p}\leq\frac{1}{2}\sqrt{\frac{1}{27}}=\frac{1}{6\sqrt{3}},$$
da cui la (\ref{eq:disug-Pr}). Inoltre si ha uguaglianza se e solo 
se $x=y=z$, cio\`{e} se e solo se $a=b=c$, cio\`{e} se e solo se il 
triangolo \`{e} equilatero.

Viceversa, consideriamo ora un triangolo equilatero circoscritto ad un 
cerchio di raggio $r$. Il perimetro di tale triangolo \`{e} 
$6\sqrt{3}\,r$. Allontanando un vertice dal cerchio, \`{e} chiaro 
che possiamo costruire triangoli isosceli con perimetro 
arbitrariamente grande, come nella seguente figura.

\begin{center}
	\psset{unit=10mm}
	\pspicture(-1,-1)(1,5)
	
	\pstGeonode[PointSymbol=none](0,0){0}(0,1){H}
	\pstCircleOA[linecolor=green]{O}{H}
	\pstGeonode[PointSymbol={*,none},PointName=none](0,2){C}(-1,-1){A'}(1,-1){B'}
	\pstHomO[HomCoef=0.5,PointSymbol=none]{O}{C}[O']
	\pstInterCC[PointSymbol=none]{O}{H}{O'}{C}{A''}{B''}
	\pstInterLL[PointName=none]{C}{A''}{A'}{B'}{A}
	\pstInterLL[PointName=none]{C}{B''}{A'}{B'}{B}
	\pstLineAB{A}{B}
	\pstLineAB{A}{C}
	\pstLineAB{B}{C}
	
	\endpspicture
	\hspace{6em}
	\psset{unit=10mm}
	\pspicture(-1,-1)(1,5)
	
	\pstGeonode[PointSymbol=none](0,0){0}(0,1){H}
	\pstCircleOA[linecolor=green]{O}{H}
	\pstGeonode[PointSymbol={*,none},PointName=none](0,3.5){C}(-1,-1){A'}(1,-1){B'}
	\pstHomO[HomCoef=0.5,PointSymbol=none]{O}{C}[O']
	\pstInterCC[PointSymbol=none]{O}{H}{O'}{C}{A''}{B''}
	\pstInterLL[PointName=none]{C}{A''}{A'}{B'}{A}
	\pstInterLL[PointName=none]{C}{B''}{A'}{B'}{B}
	\pstLineAB{A}{B}
	\pstLineAB{A}{C}
	\pstLineAB{B}{C}
	
	\endpspicture
	\hspace{6em}
	\psset{unit=10mm}
	\pspicture(-1,-1)(1,5)
	
	\pstGeonode[PointSymbol=none](0,0){0}(0,1){H}
	\pstCircleOA[linecolor=green]{O}{H}
	\pstGeonode[PointSymbol={*,none},PointName=none](0,5){C}(-1,-1){A'}(1,-1){B'}
	\pstHomO[HomCoef=0.5,PointSymbol=none]{O}{C}[O']
	\pstInterCC[PointSymbol=none]{O}{H}{O'}{C}{A''}{B''}
	\pstInterLL[PointName=none]{C}{A''}{A'}{B'}{A}
	\pstInterLL[PointName=none]{C}{B''}{A'}{B'}{B}
	\pstLineAB{A}{B}
	\pstLineAB{A}{C}
	\pstLineAB{B}{C}
	
	\endpspicture
\end{center}

Poich\'{e} il perimetro dipende in modo continuo dalla distanza dal
centro del vertice che stiamo allontanando, con un argomento di
continuit\`{a} si conclude quindi che ogni valore maggiore od uguale a
$6\sqrt{3}\,r$ \`{e} il perimetro di un triangolo isoscele
circoscritto al cerchio dato.  }

\renewcommand{\SoluzioniAlternative}{\TitoloSoluzioneAlternativaA
Siano $\alpha$, $\beta$, $x$, $y$ come nella seguente figura.
\begin{center}
	\psset{unit=10mm}
	\pspicture(0,-0.5)(4,7)
	
	\pstTriangle[PointName=none](0,0){A}(4,0){B}(2,7){C}
	\pstRotation[AngleCoef=0.5,RotAngle=\pstAngleAOB{B}{A}{C},PointSymbol=none]
	{A}{B}[A']
	\pstRotation[AngleCoef=0.5,RotAngle=\pstAngleAOB{B}{A}{C},PointSymbol=none]
	{B}{C}[B']
	\pstInterLL[PointSymbol=none]{A}{A'}{B}{B'}{I}
	\pstProjection[PointName=none]{A}{B}{I}[H]
	\pstProjection[PointName=none]{A}{C}{I}[K]
	\pstProjection[PointName=none]{C}{B}{I}[L]
	\pstCircleOA[linecolor=green]{I}{H}
	\psset{linecolor=blue}
	\pstLineAB{A}{I}
	\pstLineAB{B}{I}
	\pstLineAB{C}{I}
	\psset{linestyle=dashed}
	\pstLineAB{H}{I}
	\pstLineAB{K}{I}
	\pstLineAB{L}{I}
	
	\pstMarkAngle[MarkAngleRadius=0,LabelSep=0.7]{B}{A}{I}{$\alpha$}
	\pstMarkAngle[MarkAngleRadius=0,LabelSep=0.7]{I}{A}{C}{$\alpha$}
	\pstMarkAngle[MarkAngleRadius=0,LabelSep=0.7]{I}{B}{A}{$\alpha$}
	\pstMarkAngle[MarkAngleRadius=0,LabelSep=0.7]{C}{B}{I}{$\alpha$}
	\pstMarkAngle[MarkAngleRadius=0,LabelSep=1.3]{A}{C}{I}{$\beta$}
	\pstMarkAngle[MarkAngleRadius=0,LabelSep=1.3]{I}{C}{B}{$\beta$}
	
	\rput[t](1,-0.2){$x$}
	\rput[t](3,-0.2){$x$}
	\rput(0,1.1){$x$}
	\rput(4,1.1){$x$}
	\rput(1,4.5){$y$}
	\rput(3,4.5){$y$}
	\rput(1.3,2){$r$}
	\rput(2.7,2){$r$}
	\rput[l](2.15,0.8){$r$}
	
	\endpspicture
\end{center}

Con un po' di trigonometria elementare si ricava che
$$x=\frac{r}{\tan\alpha},
\hspace{4em}
y=\frac{r}{\tan\beta}.$$

Poich\'{e} $2\alpha+\beta=\pi/2$, usando l'opportuna formula per gli 
archi associati e la formula di duplicazione 
della tangente possiamo riscrivere $y$ nella forma
$$y=\frac{r}{\tan\beta}=
\frac{r}{\displaystyle{\tan\left(\frac{\pi}{2}-2\alpha\right)}}=
r\tan(2\alpha)=2r\cdot\frac{\tan\alpha}{1-\tan^{2}\alpha}.$$

Indicato con $P$ il perimetro del triangolo, abbiamo pertanto che
$$P=4x+2y=\frac{4r}{\tan\alpha}+4r\frac{\tan\alpha}{1-\tan^{2}\alpha}=
\frac{4r}{\tan\alpha(1-\tan^{2}\alpha)}.$$

Per determinare i possibili valori del perimetro ci siamo quindi 
ricondotti a determinare i valori assunti dalla funzione 
$\tan\alpha(1-\tan^{2}\alpha)$ quando $\alpha$ varia tra $0$ e 
$\pi/4$ (infatti l'angolo alla base del triangolo isoscele, 
$2\alpha$, deve essere acuto). Posto $z=\tan\alpha$, ci riconduciamo a 
sua volta a determinare i valori assunti dalla funzione 
$f(z)=z(1-z^{2})$ al variare di $z$ in $(0,1)$ (quando $\alpha$ varia 
in $(0,\pi/4)$, $\tan\alpha$ assume tutti e soli i valori in $(0,1)$). 

Per questo possiamo ricorrere al calcolo differenziale: dallo studio
del segno di $f'(z)=1-3z^{2}$, ricaviamo che nell'intervallo
considerato la funzione risulter\`{a} crescente in $[0,1/\sqrt{3}]$ e
decrescente in $[1/\sqrt{3},1]$.  Poich\'{e} $f(0)=f(1)=0$ e
$f(1/\sqrt{3})=2/3\sqrt{3}$, l'andamento della funzione sar\`{a} dato
dal seguente grafico.

\begin{center}
	\psset{unit=30mm}
	\pspicture(-0.4,-0.4)(1.2,1)
	
	\psline[linewidth=0.7\pslinewidth]{->}(-0.2,0)(1.2,0)
	\psline[linewidth=0.7\pslinewidth]{->}(0,-0.2)(0,1)
	\psplot[linewidth=1.7\pslinewidth]{0}{1}{x x 3 exp sub}
	\psline[linecolor=blue, linestyle=dashed](!1 3 sqrt div 0)
	(!1 3 sqrt div 2 3 3 sqrt mul div)
	\psline[linecolor=blue, linestyle=dashed](!0 2 3 3 sqrt mul div)
	(!1 3 sqrt div 2 3 3 sqrt mul div)
	\rput[r](!0.05 neg 2 3 3 sqrt mul div){$\displaystyle{\frac{2}{3\sqrt{3}}}$}
	\rput[t](!1 3 sqrt div 0.05 neg){$\displaystyle{\frac{1}{\sqrt{3}}}$}
	\rput[t](1,-0.05){1}
	\rput[tr](-0.05,-0.05){0}
	
	\endpspicture
\end{center}

Ne segue che $f(z)$ assume in $(0,1)$ tutti i valori maggiori di $0$ 
e minori od uguali di $2/3\sqrt{3}$, di conseguenza anche 
$\tan\alpha(1-\tan^{2}\alpha)$ assume gli stessi valori per 
$\alpha\in(0,\pi/4)$ e passando al reciproco si ha che $P$ assume 
tutti i valori maggiori od uguali di
$$4r\cdot \frac{3\sqrt{3}}{2}=6\sqrt{3}\,r,$$
che \`{e} quanto volevamo dimostrare.}

\renewcommand{\Osservazioni}{\TitoloOss{1} Nella seconda soluzione si
poteva dimostrare che il massimo di $f(z)$ al variare di $z$ in
$(0,1)$ si ha per $z=1/\sqrt{3}$ senza ricorrere al calcolo
differenziale.  Intanto si poteva congetturare che il
massimo si ha proprio per $z=1/\sqrt{3}$ con argomenti di tipo
geometrico: infatti \`{e} plausibile aspettarsi il minimo del
perimetro nel caso del triangolo equilatero ($\alpha=\pi/6$), dunque
il massimo di $\tan\alpha(1-\tan^{2}\alpha)$ per $\alpha=\pi/6$, e
quindi il massimo di $z(1-z^{2})$ per $z=\tan(\pi/6)=1/\sqrt{3}$.  A
questo punto non resta che dimostrare che
$$z(1-z^{2})\leq f\left(\frac{1}{\sqrt{3}}\right)=\frac{2}{3\sqrt{3}}$$
per ogni $z\in[0,1]$, e questo segue immediatamente se si osserva che, 
con qualche passaggio algebrico, tale disequazione pu\`{o} essere 
riscritta nella forma
$$(\sqrt{3}\,x-1)^{2}(\sqrt{3}\,x+2)\geq 0.$$

\TitoloOss{2}
L'approccio trigonometrico della seconda soluzione pu\`{o} anche
portare ad una dimostrazione alternativa della disuguaglianza
(\ref{eq:disug-Pr}) per un generico triangolo.  Infatti, indicati gli
angoli con $2\alpha$, $2\beta$, $2\gamma$, si ha che
\begin{equation}
	P=2r\left(\frac{1}{\tan\alpha}+\frac{1}{\tan\beta}+
	\frac{1}{\tan\gamma}\right).
	\label{eq:Ptrig}
\end{equation}

Ora la funzione $\varphi(x)=(\tan x)^{-1}$ \`{e} convessa in 
$(0,\pi/2)$, in quanto la sua derivata seconda 
$\varphi"(x)=2(\tan x)^{-3}(1+\tan^{2}x)$ \`{e} positiva in detto 
intervallo. Dalla disuguaglianza di Jensen segue pertanto che
$$\frac{1}{3}\left(\frac{1}{\tan\alpha}+\frac{1}{\tan\beta}+
\frac{1}{\tan\gamma}\right)\geq
\frac{1}{\displaystyle{\tan\left(\frac{\alpha+\beta+\gamma}{3}\right)}}=
\frac{1}{\tan(\pi/6)}=\sqrt{3},$$
che sostituita nella (\ref{eq:Ptrig}) porta alla (\ref{eq:disug-Pr}).
}


